\section{Introducción}\label{sec:introduccion}
El diseño de algoritmos tiene como objetivo encontrar soluciones a problemas de manera eficiente, considerando las limitaciones de tiempo y memoria presentes en los sistemas informáticos. En el proyecto anterior se estudiaron algoritmos iterativos de ordenamiento y búsqueda verificando experimentalmente su comportamiento cuadrático o lineal. Si bien esos resultados confirmaron la correspondencia entre el análisis teórico y el experimental, los algoritmos $O(n^2)$ presentan limitaciones evidentes frente a conjuntos de datos de gran tamaño.

El presente informe aborda el mismo problema de ordenamiento y búsqueda sobre la base de datos ficticia de deportistas, esta vez desde la perspectiva de la estrategia \textit{Divide y Vencerás}: un paradigma algorítmico que descompone el problema en subproblemas más pequeños del mismo tipo, los resuelve de forma recursiva y combina sus soluciones para obtener el resultado final. Esta estrategia da origen a algoritmos como Merge Sort y Quick Sort, que alcanzan una complejidad de $O(n \log n)$ en el caso promedio, y a variantes de búsqueda como la búsqueda exponencial y por interpolación, que extienden la búsqueda binaria para escenarios específicos.

Adicionalmente, se analiza Quick Select como algoritmo de selección del $k$-ésimo elemento, que permite resolver este problema en tiempo promedio $O(n)$ sin necesidad de ordenar el arreglo completo.